\section{Complete saturation criterion}
\label{sec:equality}

\begin{theorem}[Connected Perron saturation]\label{thm:saturation}
Assume \(G_P\) is connected.  The following are equivalent:
\begin{enumerate}[label=(\roman*),leftmargin=2.4em]
\item \(\|H\|=\rho(P)\);
\item for some \(\sigma\in\{\pm1\}\),
\(\delta_\sigma(P,H)=0\);
\item for some \(\sigma\in\{\pm1\}\) and diagonal unitary
\(D=\diag(d_i)\),
\begin{equation}\label{eq:saturation}
 |H|=P,
 \qquad
 H=\sigma DPD^*.
\end{equation}
\end{enumerate}
\end{theorem}

\begin{proof}
The equivalence of (i) and (ii) is \eqref{eq:norm-defect}.
Suppose (ii) holds.  A nonzero minimizer \(f\) in
\eqref{eq:variation} has zero energy.  If any edge has
\(q_{ij}<p_{ij}\), its amplitude term forces \(f_i=f_j=0\).
Across any adjacent saturated edge, the phase term then forces the next
vertex value to vanish; across any adjacent nonsaturated edge, the
amplitude term does the same.  Connectivity would give \(f=0\), a
contradiction.  Hence \(q_{ij}=p_{ij}\) on every edge.

The remaining zero-energy equations are
\[
 f_i=\sigma\phi_{ij}f_j.
\]
They imply equal modulus along every edge; connectivity gives
\(|f_i|=a>0\) for all \(i\).  Put \(d_i=f_i/a\).  Then
\(\phi_{ij}=\sigma d_i\overline{d_j}\), which is
\eqref{eq:saturation}.

Conversely, if \eqref{eq:saturation} holds, \(H\) is unitarily similar
to \(\sigma P\), so one of its extremal eigenvalues has modulus
\(\rho(P)\).
\end{proof}

\begin{corollary}[Pure nonnegative equality]\label{cor:nonnegative}
If \(0\le A\le P\), both matrices are symmetric, and \(G_P\) is
connected, then
\[
 \rho(A)=\rho(P)
 \quad\Longleftrightarrow\quad
 A=P.
\]
\end{corollary}

\begin{proof}
Apply \cref{thm:saturation} with \(H=A\).  Since every nonzero edge
phase equals \(1\), amplitude saturation is necessary and sufficient.
\end{proof}

\begin{corollary}[Exact amplitude--phase split]\label{cor:split}
For every Hermitian \(H\) dominated by \(P\),
\begin{equation}\label{eq:split}
 \boxed{\displaystyle
 \rho(P)-\|H\|
 =
 \bigl[\rho(P)-\rho(|H|)\bigr]
 +
 \bigl[\rho(|H|)-\|H\|\bigr].}
\end{equation}
\end{corollary}

The first bracket is completion or amplitude slack.  The second is
phase-frustration slack.  This algebraic identity is simple, but it
prevents one source from being reported as evidence for the other.
